% ============================================================
%  PE5 -- Unidad 5 | PLANTILLA DE TRABAJO DIVIDIDO
%  MundiPets -- Plataforma de Adopción y Cruza Responsable de Mascotas
%  Asignatura: Ingeniería de Requisitos (ISR-401) -- UTEQ 2026-2027 PPA
%  Formato: IEEE onecolumn | Biber + biblatex (estilo IEEE)
%  Preámbulo trasplantado y adaptado del ERS_SRS_2B_v2_0.tex del PFC,
%  para que la PE5 y el ERS compartan exactamente el mismo motor de
%  compilación (mismo compilador, mismo estilo de citas, misma
%  numeración de tablas/figuras). Ver Sección 0.6 y 0.7 del cuerpo.
%
%  COMO USAR ESTE ARCHIVO
%  ----------------------
%  Cada bloque  indica QUE debe producir un
%  integrante en esa sección. El responsable borra la caja 
%  y la reemplaza por el contenido real. Cuando no quede ninguna caja
%  gris en el documento, la PE5 está terminada.
%
%  COMPILACIÓN (requiere referencias.bib en la misma carpeta):
%      pdflatex PE5_U5_Trabajo_Dividido.tex
%      biber    PE5_U5_Trabajo_Dividido
%      pdflatex PE5_U5_Trabajo_Dividido.tex
%      pdflatex PE5_U5_Trabajo_Dividido.tex
% ============================================================
\documentclass[journal,onecolumn]{IEEEtran}

% ── Codificación e idioma ───────────────────────────────────
\usepackage[T1]{fontenc}
\usepackage[utf8]{inputenc}
\usepackage[spanish, es-tabla]{babel}
\usepackage[autostyle=true, spanish=mexican]{csquotes}
\usepackage[htt]{hyphenat}   % permite guionar dentro de \texttt (evita Overfull)

% ── Tablas ──────────────────────────────────────────────────
\usepackage{booktabs}
\usepackage{tabularx}
\usepackage{array}
\usepackage{multirow}
\usepackage{longtable}
\newcolumntype{C}[1]{>{\centering\arraybackslash}p{#1}}
\newcolumntype{L}[1]{>{\raggedright\arraybackslash}p{#1}}
\renewcommand{\thetable}{\arabic{table}}
\def\tablename{Tabla}

% ── Figuras ─────────────────────────────────────────────────
\usepackage{graphicx}
\graphicspath{{figuras/}{./}}
\usepackage{float}
\usepackage{pdflscape}   % matriz de trazabilidad en horizontal

% ── Color ───────────────────────────────────────────────────
\usepackage{xcolor}
\definecolor{uteqverde}{RGB}{0,110,60}
\definecolor{codegray}{gray}{0.95}
\definecolor{codegreen}{rgb}{0.0,0.5,0.0}
\definecolor{codeblue}{rgb}{0.0,0.0,0.6}
\definecolor{codepurple}{rgb}{0.5,0.0,0.35}

% ── Geometría ───────────────────────────────────────────────
\usepackage{geometry}
\geometry{a4paper, top=2.5cm, bottom=2.5cm, left=2.5cm, right=2.5cm,
	headheight=14pt, footskip=1.2cm}

% ── Encabezado y pie ────────────────────────────────────────
\usepackage{fancyhdr}
\pagestyle{fancy}
\fancyhf{}
\renewcommand{\headrulewidth}{0.4pt}
\fancyhead[L]{\small PE5 -- Unidad 5 -- Ingeniería de Requisitos}
\fancyhead[R]{\small MundiPets -- Trabajo dividido}
\fancyfoot[C]{\thepage}

% ── Espaciado y listas ──────────────────────────────────────
\usepackage{parskip}
\setlength{\parskip}{4pt}
\setlength{\emergencystretch}{3em}   % tolerancia extra de línea
\makeatletter
\@ifundefined{labelindent}{}{\let\labelindent\relax}
\makeatother
\usepackage{enumitem}
\usepackage{amssymb}
\usepackage{amsmath}

% ── Títulos de sección ──────────────────────────────────────
\usepackage{titlesec}
\renewcommand{\thesection}{\arabic{section}}
\renewcommand{\thesubsection}{\thesection.\arabic{subsection}}
\renewcommand{\thesubsubsection}{\thesubsection.\arabic{subsubsection}}
\titleformat{\section}{\normalfont\Large\bfseries}{\thesection}{1em}{}
\titleformat{\subsection}{\normalfont\large\bfseries}{\thesubsection}{1em}{}
\titleformat{\subsubsection}{\normalfont\normalsize\bfseries}{\thesubsubsection}{1em}{}
\setcounter{secnumdepth}{4}
\setcounter{tocdepth}{4}

% ── Leyendas ────────────────────────────────────────────────
\usepackage{caption}
\captionsetup{font=small, labelfont=bf}

% ── Bibliografía (biber + estilo IEEE) ──────────────────────
% Mismo motor que el ERS del PFC: ver Sección 0.6 y 0.7 del cuerpo
% para el orden exacto de compilación.
\usepackage[
backend=biber,
style=ieee,
citestyle=numeric-comp,
isbn=true,
doi=true
]{biblatex}
\addbibresource{referencias.bib}

% ── Enlaces ─────────────────────────────────────────────────
\usepackage[hidelinks,colorlinks=false]{hyperref}
\usepackage{url}

% ============================================================
%  CAJA DE GUÍA (fondo gris claro, borde negro)
% ============================================================
\usepackage{mdframed}
\mdfdefinestyle{guia}{
	linecolor=black,
	linewidth=0.6pt,
	backgroundcolor=gray!10,
	innertopmargin=6pt,
	innerbottommargin=6pt,
	innerleftmargin=8pt,
	innerrightmargin=8pt,
	skipabove=4pt,
	skipbelow=4pt
}

% 
\newcommand{\guia}[2]{%
	\begin{mdframed}[style=guia]
		\noindent\textbf{Responsable: #1}\\[2pt]
		\textit{#2}
	\end{mdframed}
}

% -- Atajos de nombres (evita erratas al repetirlos) ---------
\newcommand{\rFUE}{FUERTES ARRAES, Edson Daniel}
\newcommand{\rGUT}{GUTIÉRREZ ORTEGA, Génesis Adriana}
\newcommand{\rMEN}{MENDOZA PÁRRAGA, Andy Johel}
\newcommand{\rMOR}{MORALES SÁNCHEZ, Gary Alejandro}
\newcommand{\rNIE}{NIEVES SÁNCHEZ, Jimmy Samuel}
\newcommand{\rTODOS}{TODO EL EQUIPO}

% ============================================================
\begin{document}
	
	% ============================================================
	% PORTADA (mismo formato que el ERS del PFC, Apéndice de portada)
	% ============================================================

	
	% -- FT.4 --------------------------------------------------------
	\subsection*{Fundamentos de los requisitos de inteligencia artificial}
	\addcontentsline{toc}{subsection}{Fundamentos de requisitos de IA (Mendoza)}

	Especificar los componentes de inteligencia artificial (IA) de MundiPets --moderación
	de mensajes, verificación de imagen y recomendación de compatibilidad de cruza-- exige un
	marco distinto del que basta para un requisito determinista, porque el comportamiento del
	componente no está fijado en el código sino aprendido a partir de datos. Vogelsang y Borg
	documentaron este desplazamiento de paradigma, de programar a entrenar, a partir de un
	estudio con científicos de datos: sus hallazgos muestran que la ingeniería de requisitos para
	sistemas de aprendizaje automático (AA) obliga al equipo a comprender las métricas de
	rendimiento del modelo para redactar requisitos funcionales verificables, y a incorporar
	atributos de calidad nuevos --explicabilidad, ausencia de discriminación y requisitos
	legales específicos-- que no tienen equivalente directo en el desarrollo de software
	tradicional \cite{vogelsang2019}. Esta distinción es la base de la Sección~\ref{sec:requisitos_ia}: un RF-IA no describe una regla fija sino el comportamiento esperado de un modelo bajo
	incertidumbre, y un RNF-IA debe fijar umbral, unidad y método de comprobación precisamente
	porque, a diferencia de un RNF tradicional, el valor obtenido puede variar según los datos con
	los que se entrenó o se está evaluando el componente.

	Esta observación puntual se confirma a escala mediante el mapeo sistemático de Ahmad, Abdelrazek,
	Arora, Bano y Grundy, que revisó 43 estudios primarios sobre ingeniería de requisitos para
	sistemas de IA: los autores concluyen que las herramientas y plantillas de especificación
	heredadas del desarrollo de software tradicional no capturan de forma adecuada requisitos
	propios de la IA como los datos de entrenamiento, la ética o la explicabilidad, y señalan que
	la evaluación empírica de estos atributos --particularmente ética, explicabilidad y
	confianza-- sigue siendo escasa en la literatura \cite{ahmad2023}. Umm-e-Habiba, Haug, Bogner
	y Wagner llegan a una conclusión convergente desde un mapeo posterior centrado en la práctica
	profesional: la disciplina de RE para sistemas de IA es todavía inmadura frente a la RE
	tradicional, con vacíos particularmente marcados en la gestión de requisitos de datos y en la
	trazabilidad entre requisito y componente entrenado \cite{Habiba2024}. Esta brecha documentada
	desde dos ángulos distintos --el mapeo de la literatura académica y el análisis de la práctica
	profesional-- es el motivo por el cual la ficha de componente de la Sección~\ref{sec:requisitos_ia}
	obliga a completar explícitamente los campos de datos de entrenamiento, equidad y
	explicabilidad, y por el cual la Sección~\ref{sec:requisitos_ia} exige trazar cada RF-IA y
	RNF-IA hacia los requisitos preexistentes del sistema, en vez de dejarlos como una sección
	narrativa aparte del resto de la especificación.

	Un problema relacionado, y no solo de especificación sino de proceso, es el que documentan
	Amershi et al.\ a partir de un estudio de campo con equipos de Microsoft: incluso en una
	organización con procesos ágiles maduros, la integración de componentes de aprendizaje
	automático introdujo dificultades adicionales de gestión --como el acoplamiento entre datos,
	modelo y código, y la dificultad de fijar criterios de aceptación estables cuando el
	desempeño del modelo cambia con cada reentrenamiento-- que los flujos de trabajo Scrum
	convencionales no anticipaban \cite{amershi2019}. Este hallazgo justifica que el plan de
	monitoreo de la Sección~\ref{sec:requisitos_ia} no se trate como una tarea de una sola vez,
	sino como una traza continua que debe revisarse en cada iteración del componente.

	Sobre esa base, Habibullah, Gay y Horkoff aportan evidencia empírica de que los requisitos no
	funcionales de aprendizaje automático reciben en la práctica profesional un tratamiento
	heterogéneo: los equipos rara vez documentan explícitamente atributos como equidad,
	explicabilidad o robustez frente a datos fuera de distribución, y cuando lo hacen carecen de
	un método sistemático de medición \cite{habibullah2023}. Este hallazgo es la justificación
	directa de exigir, en la Sección~\ref{sec:requisitos_ia}, que cada RNF-IA de MundiPets declare
	umbral numérico, unidad y método de comprobación explícitos: sin esa disciplina, la moderación
	de mensajes o la verificación de imagen quedarían documentadas como intención (``el modelo debe
	ser preciso'') y no como un requisito auditable.

	El componente legal del marco proviene de dos fuentes complementarias. El Reglamento (UE)
	2024/1689, conocido como Ley de IA, adopta un enfoque basado en niveles de riesgo y establece
	obligaciones de transparencia, supervisión humana y documentación técnica proporcionales al
	riesgo que introduce cada sistema de IA \cite{euaiact2024}; aunque MundiPets opera en Ecuador y
	no está sujeto directamente a su jurisdicción, el equipo lo adopta como referencia de buenas
	prácticas para clasificar el riesgo de sus tres componentes y fijar las obligaciones de
	supervisión humana descritas en la Sección~\ref{sec:etica_ia}. La base legal aplicable dentro
	de Ecuador es la Ley Orgánica de Protección de Datos Personales, vigente desde su publicación
	en el Registro Oficial en 2021, que exige finalidad determinada, minimización de datos y
	consentimiento del titular para el tratamiento de datos personales \cite{lopdp2021}; esa misma
	ley fue desarrollada operativamente dos años después por su Reglamento General, que precisa los
	mecanismos concretos de seudonimización, plazos y conservación de la información que la ley de
	2021 dejó enunciados en términos generales \cite{lopdp_reglamento2023}. Se citan ambas normas,
	y no solo la más reciente, porque no son versiones sucesivas de un mismo texto sino dos niveles
	distintos y vigentes de la misma obligación: la ley fija qué se debe garantizar y el reglamento
	fija cómo se implementa; omitir la ley de 2021 dejaría sin base normativa directa a la
	obligación misma que el reglamento de 2023 desarrolla.

	En conjunto, estas fuentes explican por qué un componente probabilístico exige medidas que un
	requisito determinista no necesita: un RF tradicional de MundiPets, como registrar el perfil
	médico de una mascota, se comprueba una sola vez contra una regla fija y no cambia de
	comportamiento si cambian los datos de entrada. Un componente de IA, en cambio, puede degradar
	su desempeño con el tiempo (deriva de datos), producir resultados distintos para grupos de
	usuarios distintos (inequidad) y tomar decisiones cuyo razonamiento interno no es directamente
	observable (opacidad); por eso la Sección~\ref{sec:requisitos_ia} exige, además de los RF y RNF
	de rendimiento habituales, RNF de equidad, un RNF de explicabilidad y un plan de monitoreo en
	producción que ningún requisito determinista de MundiPets necesita.
	

	%  REFERENCIAS
	% ============================================================
	\printbibliography[title={Referencias}]
	

	
\end{document}